\documentclass[10pt,letterpaper]{article}
\usepackage[utf8]{inputenc}
\usepackage[intlimits]{amsmath}
\usepackage{amsfonts}
\usepackage{amssymb}
\usepackage{ragged2e}
\usepackage[letterpaper, margin=0.5in]{geometry}
\usepackage{graphicx}
\usepackage{cancel}
\usepackage{mathtools}
\usepackage{tabularx}
\usepackage{arydshln}
\usepackage{tensor}
\usepackage{array}
\usepackage{xcolor}
\usepackage[boxed]{algorithm}
\usepackage[noend]{algpseudocode}
\usepackage{listings}
\usepackage{textcomp}
% \usepackage[pdf,tmpdir,singlefile]{graphviz}
\usepackage{mathrsfs}
\usepackage{bbm}
\usepackage{tikz}
\usepackage{tikz-cd}
\usepackage{enumitem}
\usepackage{arydshln}
\usepackage{relsize}
\usepackage{multicol}
\usepackage{scalerel}
\usepackage{upgreek}
\usepackage{ifthen}

\usetikzlibrary{bayesnet}
\setlist{noitemsep}

%%%%%%%%%%%%%%%%%%%%%%%%%%%%%
% Formatting commands
%%%%%%%%%%%%%%%%%%%%%%%%%%%%%
\newcommand{\n}{\hfill\break}
\newcommand{\nn}{\vspace{0.5\baselineskip}\n}
\newcommand{\up}{\vspace{-\baselineskip}}
\newcommand{\hangblock}[2]{\par\noindent\settowidth{\hangindent}{\textbf{#1: }}\textbf{#1: }\nolinebreak#2}
\newcommand{\lemma}[1]{\hangblock{Lemma}{#1}}
\newcommand{\defn}[1]{\hangblock{Defn}{#1}}
\newcommand{\thm}[1]{\hangblock{Thm}{#1}}
\newcommand{\cor}[1]{\hangblock{Cor}{#1}}
\newcommand{\prop}[1]{\hangblock{Prop}{#1}}
\newcommand{\ex}[1]{\hangblock{Ex}{#1}}
\newcommand{\exer}[1]{\hangblock{Exer}{#1}}
\newcommand{\fact}[1]{\hangblock{Fact}{#1}}
\newcommand{\remark}[1]{\hangblock{Remark}{#1}}
\newcommand{\proven}{\;$\square$\n}
\newcommand{\problem}[1]{\par\noindent{\nolinebreak#1}\n}
\newcommand{\problempart}[2]{\par\noindent\indent{}\settowidth{\hangindent}{\textbf{(#1)} \indent{}}\textbf{(#1) }\nolinebreak#2\n}
\newcommand{\ptxt}[1]{\textrm{\textnormal{#1}}}
\newcommand{\inlineeq}[1]{\centerline{$\displaystyle #1$}}
\newcommand{\pageline}{\up\par\noindent\rule{\textwidth}{0.1pt}}

%%%%%%%%%%%%%%%%%%%%%%%%%%%%%
% Math commands
%%%%%%%%%%%%%%%%%%%%%%%%%%%%%
% Set Theory
\newcommand{\card}[1]{\left|#1\right|}
\newcommand{\set}[1]{\left\{#1\right\}}
\newcommand{\setmid}{\;\middle|\;}
\newcommand{\ps}[1]{\mathcal{P}\left(#1\right)}
\newcommand{\pfinite}[1]{\mathcal{P}^{\ptxt{finite}}\left(#1\right)}
\newcommand{\naturals}{\mathbb{N}}
\newcommand{\N}{\naturals}
\newcommand{\integers}{\mathbb{Z}}
\newcommand{\Z}{\integers}
\newcommand{\rationals}{\mathbb{Q}}
\newcommand{\Q}{\rationals}
\newcommand{\reals}{\mathbb{R}}
\newcommand{\R}{\reals}
\newcommand{\complex}{\mathbb{C}}
\newcommand{\C}{\complex}
\newcommand{\halfPlane}{\mathbb{H}}
\let\HSym\H
\let\H\relax
\newcommand{\H}{\halfPlane}
\newcommand{\comp}{^{\complement}}
\DeclareMathOperator{\Hom}{Hom}
\newcommand{\Ind}{\mathbbm{1}}
\newcommand{\cut}{\setminus}
\DeclareMathOperator{\elem}{elem}

% Differentiable Manifolds
\newcommand{\RP}{\mathbb{RP}}
\newcommand{\CP}{\mathbb{RP}}
\newcommand{\osubset}{\overset{\mathclap{\scalebox{0.5}{\ptxt{open}}}}{\subset}}
\newcommand{\osubseteq}{\overset{\mathclap{\scalebox{0.5}{\ptxt{open}}}}{\subseteq}}
\newcommand{\osupset}{\overset{\mathclap{\scalebox{0.5}{\ptxt{open}}}}{\supset}}
\newcommand{\osupseteq}{\overset{\mathclap{\scalebox{0.5}{\ptxt{open}}}}{\supseteq}}
\newcommand{\pdat}[3]{\left.\pd{#1}{#2}\right|_{#3}}
\DeclareMathOperator{\so}{so}
\DeclareMathOperator{\codim}{codim}
\DeclareMathOperator{\Diff}{Diff}
\let\dSym\d
\let\d\relax
\newcommand{\d}{\partial}
\DeclareMathOperator{\gl}{gl}

% Graph Theory
\let\deg\relax
\DeclareMathOperator{\deg}{deg}
\newcommand{\degp}{\ptxt{deg}^{+}}
\newcommand{\degn}{\ptxt{deg}^{-}}
\newcommand{\precdot}{\mathrel{\ooalign{$\prec$\cr\hidewidth\hbox{$\cdot\mkern0.5mu$}\cr}}}
\newcommand{\succdot}{\mathrel{\ooalign{$\cdot\mkern0.5mu$\cr\hidewidth\hbox{$\succ$}\cr\phantom{$\succ$}}}}
\DeclareMathOperator{\cl}{cl}
\DeclareMathOperator{\affdim}{affdim}

% Probability
\newcommand{\parSymbol}{\P}
\newcommand{\Prob}{\mathbb{P}}
\renewcommand{\P}{\Prob}
\newcommand{\Avg}{\mathbb{E}}
\newcommand{\E}{\Avg}
\DeclareMathOperator{\Var}{Var}
\DeclareMathOperator{\cov}{cov}
\DeclareMathOperator{\Unif}{Unif}
\DeclareMathOperator{\Binom}{Binom}
\newcommand{\CI}{\mathrel{\text{\scalebox{1.07}{$\perp\mkern-10mu\perp$}}}}

% Standard Math
\newcommand{\inv}{^{-1}}
\newcommand{\abs}[1]{\left|#1\right|}
\newcommand{\ceil}[1]{\left\lceil{}#1\right\rceil{}}
\newcommand{\floor}[1]{\left\lfloor{}#1\right\rfloor{}}
\newcommand{\conj}[1]{\overline{#1}}
\newcommand{\of}{\circ}
\newcommand{\tri}{\triangle}
\newcommand{\inj}{\hookrightarrow}
\newcommand{\surj}{\twoheadrightarrow}
\newcommand{\ndiv}{\nmid}
\renewcommand{\epsilon}{\varepsilon}
\newcommand{\divides}{\mid}
\newcommand{\ndivides}{\nmid}
\DeclareMathOperator{\lcm}{lcm}
\DeclareMathOperator{\sgn}{sgn}
\newcommand{\map}[4]{\!\!\!\begin{array}[t]{rcl}#1 & \!\!\!\!\to & \!\!\!\!#2\\ {}#3 & \!\!\!\!\mapsto & \!\!\!\!#4\end{array}}
\newcommand{\bigsum}[2]{\smashoperator[lr]{\sum_{\scalebox{#1}{$#2$}}}}
\DeclareMathOperator{\gcf}{gcf}
\newcommand{\restr}[2]{\left.#1\right|_{#2}}

% Linear Algebra
\newcommand{\Id}{\textrm{\textnormal{Id}}}
\newcommand{\im}{\textrm{\textnormal{im}}}
\newcommand{\norm}[1]{\abs{\abs{#1}}}
\newcommand{\tpose}{^{T}\!}
\newcommand{\iprod}[1]{\left<#1\right>}
\newcommand{\giprod}{\iprod{\;\,,\;}}
\DeclareMathOperator{\tr}{tr}
\DeclareMathOperator{\trace}{tr}
\newcommand{\chgBasMat}[3]{\!\!\tensor*[_{#1}]{\left[#2\right]}{_{#3}}}
\newcommand{\vecBas}[2]{\tensor*[]{\left[#1\right]}{_{#2}}}
\DeclareMathOperator{\GL}{GL}
\DeclareMathOperator{\Mat}{Mat}
\DeclareMathOperator{\vspan}{span}
\DeclareMathOperator{\rank}{rank}
\newcommand{\V}[1]{\vec{#1}}
\DeclareMathOperator{\proj}{proj}
\DeclareMathOperator{\compProj}{comp}
\DeclareMathOperator{\row}{row}
\newcommand{\smallPMatrix}[1]{\paren{\begin{smallmatrix}#1\end{smallmatrix}}}
\newcommand{\smallBMatrix}[1]{\brack{\begin{smallmatrix}#1\end{smallmatrix}}}
\newcommand{\pmat}[1]{\begin{pmatrix}#1\end{pmatrix}}
\newcommand{\bmat}[1]{\begin{bmatrix}#1\end{bmatrix}}
\newcommand{\dual}{^{*}}
\newcommand{\pinv}{^{\dagger}}
\newcommand{\horizontalMatrixLine}{\ptxt{\rotatebox[origin=c]{-90}{$|$}}}
\DeclareMathOperator{\range}{range}
\DeclareMathOperator{\Symm}{Symm}
\DeclareMathOperator{\SU}{SU}

% Multilinear Algebra
\let\Lsym\L
\let\L\relax
\DeclareMathOperator{\L}{\mathscr{L}}
\DeclareMathOperator{\A}{\mathcal{A}}
\DeclareMathOperator{\Alt}{Alt}
\DeclareMathOperator{\Sym}{Sym}
\newcommand{\ot}{\otimes}
\newcommand{\ox}{\otimes}
\DeclareMathOperator{\asc}{asc}
\DeclareMathOperator{\asSet}{set}
\DeclareMathOperator{\sort}{sort}
\DeclareMathOperator{\ringA}{\mathring{A}}
\DeclareMathOperator{\Sh}{Sh}
\DeclareMathOperator{\Bil}{Bil}

% Topology
\newcommand{\closure}[1]{\overline{#1}}
\newcommand{\uball}{\mathcal{U}}
\DeclareMathOperator{\Int}{Int}
\DeclareMathOperator{\Ext}{Ext}
\DeclareMathOperator{\Bd}{Bd}
\DeclareMathOperator{\rInt}{rInt}
\DeclareMathOperator{\ch}{ch}
\DeclareMathOperator{\ah}{ah}
\newcommand{\LargerTau}{\mathlarger{\mathlarger{\mathlarger{\mathlarger{\tau}}}}}
\newcommand{\Tau}{\mathcal{T}}

% Analysis
\DeclareMathOperator{\Graph}{Graph}
\DeclareMathOperator{\epi}{epi}
\DeclareMathOperator{\hypo}{hypo}
\DeclareMathOperator{\supp}{supp}
\newcommand{\lint}[2]{\underset{#1}{\overset{#2}{{\color{black}\underline{{\color{white}\overline{{\color{black}\int}}\color{black}}}}}}}
\newcommand{\uint}[2]{\underset{#1}{\overset{#2}{{\color{white}\underline{{\color{black}\overline{{\color{black}\int}}\color{black}}}}}}}
\newcommand{\alignint}[2]{\underset{#1}{\overset{#2}{{\color{white}\underline{{\color{white}\overline{{\color{black}\int}}\color{black}}}}}}}
\newcommand{\extint}{\ptxt{ext}\int}
\newcommand{\extalignint}[2]{\ptxt{ext}\alignint{#1}{#2}}
\newcommand{\conv}{\ast}
\newcommand{\pd}[2]{\frac{\partial{}#1}{\partial{}#2}}
\newcommand{\del}{\nabla}
\DeclareMathOperator{\grad}{grad}
\DeclareMathOperator{\curl}{curl}
\let\div\relax
\DeclareMathOperator{\div}{div}
\DeclareMathOperator{\vol}{vol}

% Complex Analysis
\let\Re\relax
\DeclareMathOperator{\Re}{Re}
\let\Im\relax
\DeclareMathOperator{\Im}{Im}
\DeclareMathOperator{\Res}{Res}

% Abstract Algebra
\DeclareMathOperator{\ord}{ord}
\newcommand{\generated}[1]{\left<#1\right>}
\newcommand{\cycle}[1]{\smallPMatrix{#1}}
\newcommand{\id}{\ptxt{id}}
\newcommand{\iso}{\cong}
\DeclareMathOperator{\Aut}{Aut}
\DeclareMathOperator{\SL}{SL}
\DeclareMathOperator{\op}{op}
\newcommand{\isom}[4]{\!\!\!\begin{array}[t]{rcl}#1 & \!\!\!\!\overset{\sim}{\to} & \!\!\!\!#2\\ #3 & \!\!\!\!\mapsto & \!\!\!\!#4\end{array}}
\newcommand{\F}{\mathbb{F}}
\newcommand{\acts}{\;\reflectbox{\rotatebox[origin=c]{-90}{$\circlearrowright$}}\;}
\newcommand{\disjunion}{\mathrel{\text{\scalebox{1.07}{$\perp\mkern-10mu\perp$}}}}
\DeclareMathOperator{\SO}{SO}
\DeclareMathOperator{\stab}{stab}
\DeclareMathOperator{\nullity}{nullity}
\DeclareMathOperator{\Perm}{Perm}
\DeclareMathOperator{\nsubgp}{\vartriangleleft}
\DeclareMathOperator{\notnsubgp}{\ntriangleleft}
\newcommand{\presentation}[2]{\left<#1\;\middle|\;#2\right>}
\DeclareMathOperator{\Char}{char}
\DeclareMathOperator{\fchar}{char}
\DeclareMathOperator{\triv}{triv}
\DeclareMathOperator{\reg}{reg}
\DeclareMathOperator{\std}{std}
\DeclareMathOperator{\Func}{Func}
\DeclareMathOperator{\End}{End}

% Convex Optimization
\let\sectionSymbol\S
\let\S\relax
\newcommand{\S}{\mathbb{S}}
\DeclareMathOperator{\dist}{dist}
\DeclareMathOperator{\dom}{dom}
\DeclareMathOperator{\diag}{diag}
\DeclareMathOperator{\ones}{\mathbbm{1}}
\newcommand{\minimizeOver}[3]{\begin{array}{rl}\underset{#1}{\ptxt{minimize}} & #2\\ \ptxt{subject to} & #3\end{array}}
\newcommand{\maximizeOver}[3]{\begin{array}{rl}\underset{#1}{\ptxt{maximize}} & #2\\ \ptxt{subject to} & #3\end{array}}
\newcommand{\minimizationProblem}[2]{\minimizeOver{}{#1}{#2}}
\newcommand{\maximizationProblem}[2]{\maximizeOver{}{#1}{#2}}
\newcommand{\minimizeOverUnconstrained}[2]{\begin{array}{rl}\underset{#1}{\ptxt{minimize}} & #2\end{array}}
\newcommand{\maximizeOverUnconstrained}[2]{\begin{array}{rl}\underset{#1}{\ptxt{maximize}} & #2\end{array}}
\newcommand{\minimizationUnconstrained}[1]{\minimizeOverUnconstrained{}{#1}}
\newcommand{\maximizationUnconstrained}[1]{\maximizeOverUnconstrained{}{#1}}
\DeclareMathOperator{\argmin}{argmin}
\DeclareMathOperator{\argmax}{argmax}

% Proofs
\newcommand{\st}{s.t.}
\newcommand{\unique}{!}
\newcommand{\iffdef}{\overset{\ptxt{def}}{\Leftrightarrow}}
\newcommand{\eqVertical}{\rotatebox[origin=c]{90}{=}}
\newcommand{\mapsfrom}{\mathrel{\reflectbox{\ensuremath{\mapsto}}}}
\newcommand{\mapsdown}{\rotatebox[origin=c]{-90}{$\mapsto$}\mkern2mu}
\newcommand{\mapsup}{\rotatebox[origin=c]{90}{$\mapsto$}\mkern2mu}
\newcommand{\from}{\!\mathrel{\reflectbox{\ensuremath{\to}}}}
\newcommand{\labeledeq}[1]{\overset{\mathclap{\ptxt{#1}}}{=}}
\newcommand{\eqdef}{\labeledeq{def}}

% Brackets
\newcommand{\paren}[1]{\left(#1\right)}
\renewcommand{\brack}[1]{\left[#1\right]}
\renewcommand{\brace}[1]{\left\{#1\right\}}
\newcommand{\ang}[1]{\left<#1\right>}

% Algorithms
\algrenewcommand{\algorithmiccomment}[1]{\hskip 1em \texttt{// #1}}
\algrenewcommand\algorithmicrequire{\textbf{Input:}}
\algrenewcommand\algorithmicensure{\textbf{Output:}}
\newcommand{\algP}{\ptxt{\textbf{P}}}
\newcommand{\algNP}{\ptxt{\textbf{NP}}}
\newcommand{\algNPC}{\ptxt{\textbf{NP-Complete}}}
\newcommand{\algNPH}{\ptxt{\textbf{NP-Hard}}}
\newcommand{\algEXP}{\ptxt{\textbf{EXP}}}
\DeclareMathOperator{\fl}{fl}

%%%%%%%%%%%%%%%%%%%%%%%%%%%%%
% Other commands
%%%%%%%%%%%%%%%%%%%%%%%%%%%%%
\newcommand{\flag}[1]{\textbf{\textcolor{red}{#1}}}
\let\uSym\u
\let\u\relax
\newcommand{\u}[1]{\underline{#1}}
\let\bSym\b
\let\b\relax
\newcommand{\b}[1]{\textbf{#1}}
\let\iSym\i
\let\i\relax
\newcommand{\i}[1]{\textit{#1}}
\let\scSym\sc
\let\sc\relax
\newcommand{\sc}[1]{\textsc{#1}}

%%%%%%%%%%%%%%%%%%%%%%%%%%%%%%%%%%%%%%%
% Make l's curvy in math environments %
%%%%%%%%%%%%%%%%%%%%%%%%%%%%%%%%%%%%%%%
\mathcode`l="8000
\begingroup
\makeatletter
\lccode`\~=`\l
\DeclareMathSymbol{\lsb@l}{\mathalpha}{letters}{`l}
\lowercase{\gdef~{\ifnum\the\mathgroup=\m@ne \ell \else \lsb@l \fi}}%
\endgroup

%%%%%%%%%%%%%%%%%%%%%%%%%
% Fix \vdots and \ddots %
%%%%%%%%%%%%%%%%%%%%%%%%%
\usepackage{letltxmacro}
\LetLtxMacro\orgvdots\vdots
\LetLtxMacro\orgddots\ddots

\makeatletter
\DeclareRobustCommand\vdots{%
	\mathpalette\@vdots{}%
}
\newcommand*{\@vdots}[2]{%
	% #1: math style
	% #2: unused
	\sbox0{$#1\cdotp\cdotp\cdotp\m@th$}%
	\sbox2{$#1.\m@th$}%
	\vbox{%
		\dimen@=\wd0 %
		\advance\dimen@ -3\ht2 %
		\kern.5\dimen@
		% remove side bearings
		\dimen@=\wd2 %
		\advance\dimen@ -\ht2 %
		\dimen2=\wd0 %
		\advance\dimen2 -\dimen@
		\vbox to \dimen2{%
			\offinterlineskip
			\copy2 \vfill\copy2 \vfill\copy2 %
		}%
	}%
}
\DeclareRobustCommand\ddots{%
	\mathinner{%
		\mathpalette\@ddots{}%
		\mkern\thinmuskip
	}%
}
\newcommand*{\@ddots}[2]{%
	% #1: math style
	% #2: unused
	\sbox0{$#1\cdotp\cdotp\cdotp\m@th$}%
	\sbox2{$#1.\m@th$}%
	\vbox{%
		\dimen@=\wd0 %
		\advance\dimen@ -3\ht2 %
		\kern.5\dimen@
		% remove side bearings
		\dimen@=\wd2 %
		\advance\dimen@ -\ht2 %
		\dimen2=\wd0 %
		\advance\dimen2 -\dimen@
		\vbox to \dimen2{%
			\offinterlineskip
			\hbox{$#1\mathpunct{.}\m@th$}%
			\vfill
			\hbox{$#1\mathpunct{\kern\wd2}\mathpunct{.}\m@th$}%
			\vfill
			\hbox{$#1\mathpunct{\kern\wd2}\mathpunct{\kern\wd2}\mathpunct{.}\m@th$}%
		}%
	}%
}
\makeatother

\newcommand{\B}{
	\begin{tikzpicture}
	\filldraw [fill=red, draw=black] (0, 0) rectangle (0.37, 0.45);
	\draw [line width=0.5mm, white ] (0.1,0.08) -- (0.1,0.38);
	\draw[line width=0.5mm, white ] (0.1, 0.35) .. controls (0.2, 0.35) and (0.4, 0.2625) .. (0.1, 0.225);
	\draw[line width=0.5mm, white ] (0.1, 0.225) .. controls (0.2, 0.225) and (0.4, 0.1625) .. (0.1, 0.1);
	\end{tikzpicture}
}

% Allow custom symbols as arrows in tikz-cd
\tikzset{
  symbol/.style={
    draw=none,
    every to/.append style={
      edge node={node [sloped, allow upside down, auto=false]{$#1$}}}
  }
}
% Example: \arrow[r,symbol=\cong]
% https://tex.stackexchange.com/questions/394154/how-to-include-inclusion-subgroup-relationship-in-tikz-cd-diagram

\author{Thomas Cohn}
\title{Math 635 Lecture 3}
\date{1/25/21} % Can also use \today

\begin{document}
\maketitle
\setlength\RaggedRightParindent{\parindent}
\RaggedRight

\par\noindent
Review from Friday: Given a finite-dimensional vector space $V$ (i.e. $T_{p}M$ for some $p\in{}M$), a tensor of type $(k,l)$ is an element of
\[
	T^{(k,l)}(V)\eqdef\underbrace{V\otimes\cdots\otimes{}V}_{k}\otimes\underbrace{V^{*}\otimes\cdots\otimes{}V^{*}}_{l}\cong\{\ptxt{multilinear maps }\underbrace{V^{*}\times\cdots\times{}V^{*}}_{k}\times\underbrace{V\times\cdots\times{}V}_{l}\to\R\}
\]

\ex{
	\begin{itemize}[topsep=0pt, itemsep=0pt, leftmargin=3\parindent]
		\item Type $(0,1)$ is $V^{*}$, the set of linear maps $V\to\R$.
		\item Type $(1,0)$ is $(V^{*})^{*}\cong{}V$, by $\displaystyle\isom{V}{(V^{*})}{v}{(\alpha\mapsto\alpha(v))}$ (for $\alpha\in{}V^{*}$).
		\item Type $(0,2)$ is $V^{*}\otimes{}V^{*}\cong\set{\ptxt{bilinear }V\times{}V\to\R}$.
	\end{itemize}\up\n
}

\defn{
	We can take the \u{tensor product} (or \u{outer product}) of two tensors on the same vector space: given $\tau\in{}T^{(k,l)}(V)$ and $\tau'\in{}T^{(k',l')}(V)$, we define $\tau\otimes\tau'\in{}T^{(k+k',l+l')}(V)$ by
	\[
		(\tau\otimes\tau')(\underbrace{v_{1},\ldots,v_{k+k'}}_{\in{}V},\underbrace{\alpha_{1},\cdots,\alpha_{l+l'}}_{\in{}V^{*}})\eqdef\tau(v_{1},\ldots,v_{k},\alpha_{1},\ldots,\alpha_{k})\cdot\tau'(v_{k+1},\ldots,v_{k+k'},\alpha_{l+1},\ldots,\alpha_{l+l'})
	\]
	\n
}

\par\noindent
Note: We have to check that this map is indeed multilinear.\n

\ex{
	If $\alpha,\beta\in{}T^{(0,1)}(V)=V^{*}$, then $\alpha\otimes\beta$ is the map $\map{V\times{}V}{\R}{(v_{1},v_{2})}{\alpha(v_{1})\beta(v_{2})}$.\n
}

\par\noindent
Now, we want to consider tensors on manifolds. By the lemma from last time, given $k,l$, there's a bundle\break $\pi:T^{(k,l)}(TM)\to{}M$ with fibers $T^{(k,l)}(T_{p}M)=T^{(k,l)}(\pi\inv(p))$ for $p\in{}M$. Given coordinates $(x^{1},\ldots,x^{n})$ on $U\subseteq{}M$, we have a smooth moving frame:
\[
	\set{\pd{}{x^{i_{1}}}\otimes\cdots\otimes\pd{}{x^{i_{k}}}\otimes{}dx^{j_{1}}\otimes\cdots\otimes{}dx^{j_{l}}\setmid{}i_{a},j_{b}\in\set{1,\ldots,n}}
\]
Last time, we stated that this is a basis, and any smooth section is a linear combination of these components, with $C^{\infty}$ functions as coefficients:
\[
	\sum{}A^{i_{1}\cdots{}i_{k}}_{j_{1}\cdots{}j_{k}}\pd{}{x^{i_{1}}}\otimes\cdots\otimes\pd{}{x^{i_{k}}}\otimes{}dx^{j_{1}}\otimes\cdots\otimes{}dx^{j_{l}}
\]
Why is this true?\n

\exer{
	Check that $A^{i_{1}\cdots{}i_{k}}_{j_{1}\cdots{}j_{l}}=\tau(dx^{i_{1}},\ldots,dx^{i_{k}},\pd{}{x^{j_{1}}},\ldots,pd{}{x^{j_{l}}})$.\n
}

\ex{
	Fix a smooth $1$-form $\alpha\in\Omega^{1}(U)$. Then $\alpha=\sum_{j=1}^{n}a_{j}dx^{j}$, and $a_{j}=\alpha(\pd{}{x^{j}})$.\n
}

\defn{
	Some terminology: We say $\tau\in\Gamma(T^{(k,l)}(TM))$ is a \u{tensor}, or \u{tensor field}, on $M$ of type $(k,l)$. $k$ is the \u{contravariant degree}, and $l$ is the \u{covariant degree}.\n
}

\ex{
	Vector fields are contravariant (type $(1,0)$).\n
	$1$-forms are covariant (type $(0,1)$).\n
}

\subsection*{Einstein Summation Notation}

\par\noindent
In coordinates, the covariant indices are subindices, and the contravariant indices are superindices. The convention of Einstein summation notation is if the same index appears exactly twice in a monomial -- one upper and one lower -- then summation over that index is implied.\n

\ex{
	(Linear case) $V=\R^{n}$, $\alpha\in{}V^{*}\otimes{}V\cong\Hom(V,V)$. Then we have
	\[
		a_{j}^{i}v^{j}=\sum_{j=1}^{n}a_{j}^{i}v^{j}=\brack{\begin{pmatrix}
			a_{1}^{1} & \cdots & a_{n}^{1}\\
			\vdots & & \vdots\\
			a_{1}^{n} & \cdots & a_{n}^{n}
		\end{pmatrix}\begin{pmatrix}
			v^{1}\\
			\vdots\\
			v^{n}
		\end{pmatrix}}_{i}
	\]
	Note that $(a_{j}^{i})$ is taken to be a matrix, where $j$ is the column index and $i$ the row index.\n
}

\par\noindent
Looking ahead, we'll also have a contraction operation: $a\in{}V^{*}\otimes{}V\rightsquigarrow\alpha\in\Hom(V,V)\rightsquigarrow\tr(a)=a_{i}^{i}\in\R$. \i{This doesn't require coordinates, as trace is basis-invariant.}\n

\defn{
	A \u{Riemannian metric} $g$ on a manifold $M$ is a smooth $(0,2)$ tensor, which, at every point, as a matrix, is symmetric ($g_{p}=g_{p}\tpose$) and positive definite ($\forall{}v\in{}T_{p}M$, $v\tpose{}g_{p}v\ge{}0$, with equality iff $v=0$).\n
}

\par\noindent
Observe: $\forall{}p\in{}M$, $g_{p}:T_{p}M\times{}T_{p}M\to\R$ is a bilinear map. In coordinates $(x^{1},\ldots,x^{n})$ on $U\osubseteq{}M$, $g=g_{ij}\,dx^{i}\otimes{}dx^{j}$, where $g_{ij}\in{}C^{\infty}(U)$, and $g_{ij}=g_{ji}$. In fact, because of symmetry, we can write (using Einstein summation notation)
\begin{align*}
	g & =\frac{1}{2}\brack{g_{ij}\,dx^{i}\otimes{}dx^{j}+g_{ji}\,dx^{i}\otimes{}dx^{j}}\\
	& =\frac{1}{2}\brack{g_{ij}\,dx^{i}\otimes{}dx^{j}+g_{ij}\,dx^{j}\otimes{}dx^{i}}\\
	& =g_{ij}\frac{1}{2}\underbrace{\brack{dx^{i}\otimes{}dx^{j}+dx^{j}\otimes{}dx^{i}}}_{\hspace{-100pt}\mathrlap{\eqdef{}dx^{i}dx^{j}=dx^{j}dx^{i},\ptxt{ the symmetric product}}}\\
\end{align*}
It's also standard notation to write $ds^{2}\eqdef{}dx^{i}dx^{j}$, so we can write $g=g_{ij}\,ds^{2}$.\n

\defn{
	A \u{Riemannian manifold} is a pair $(M,g)$, where $M$ is a smooth manifold, and $g$ is a Riemannian metric on $M$.\n
}

\defn{
	If $(M,g^{M})$ and $(N,g^{N})$ are Riemannian manifolds, an \u{isometry} $F:M\to{}N$ is a diffeomorphism \st{} $\forall{}p\in{}M$, $(dF)_{p}:T_{p}M\to{}T_{F(p)}N$ preserves the metric, i.e., \st{} $\forall{}u,v\in{}T_{p}M$, $g_{F(p)}^{N}(dF_{p}(u),dF_{p}(v))=g_{p}^{M}(u,v)$.\n
}

\par\noindent
Notation: Sometimes, we'll write $g_{p}(u,v)=\iprod{v,w}_{p}=\iprod{v,w}$. (We omit the subscript when the choice of metric is obvious.)\n

\ex{
	Let $M\subset\R^{N}$, a regular submanifold. $\forall{}p\in{}M$, $T_{p}M\subset\R^{N}$ can be identified with a subspace of $T_{p}R^{N}\cong\R^{N}$. Define $g_{p}=\giprod{}$ by the usual inner product in $\R^{N}$.\n
}

\ex{
	Let $M=S^{2}\subseteq\R^{3}$. We'll compute in coordinates in the upper hemisphere. Say our coordinates are $(s,t)$ (with projection to the flat unit disk), then we have our inverse parameterization
	\[
		(s,t)\mapsto(s,t,\sqrt{1-s^{2}-t^{2}})\eqdef\vec{r}(s,t)
	\]
	(with notation as in Calc III). Then under $T_{p}S^{2}\hookrightarrow\R^{3}$, we have
	\begin{align*}
		\pd{}{s} & \mapsto\pd{\vec{r}}{s}=\paren{1,0,\frac{-s}{\sqrt{1-s^{2}-t^{2}}}}\\
		\pd{}{t} & \mapsto\pd{\vec{r}}{t}=\paren{0,1,\frac{-t}{\sqrt{1-s^{2}-t^{2}}}}
	\end{align*}
	So our metric is
	\[
		(g_{ij})=\begin{pmatrix}
			\pd{\vec{r}}{s}\cdot\pd{\vec{r}}{s} & \pd{\vec{r}}{s}\cdot\pd{\vec{r}}{t}\\
			\pd{\vec{r}}{t}\cdot\pd{\vec{r}}{s} & \pd{\vec{r}}{t}\cdot\pd{\vec{r}}{t}
		\end{pmatrix}=\begin{pmatrix}
			1-t^{2} & \frac{st}{1-s^{2}-t^{2}}\\
			\frac{st}{1-s^{2}-t^{2}} & 1-s^{2}
		\end{pmatrix}
	\]
	\n
}

\par\noindent
Observe: Any submanifold $S\subset{}M$ of a Riemannian manifold $(M,g)$ inherits a Riemannian metric by $\forall{}p\in{}S$, $T_{p}S\hookrightarrow{}T_{p}M$.\n

\end{document}